\topic{2}{What the finite-volume method conserves}
A \emph{mesh} divides the fluid domain into small regions called cells.
A \emph{control volume} is a fixed region over which we impose a conservation
balance. In this two-dimensional lesson it is a cell area with one unit of
depth perpendicular to the page. A \emph{face} is one boundary segment of
that cell. Two neighboring cells have one face in common.

The finite-volume method starts from the conservation law for each control
volume. It accounts for what is stored inside the region, what crosses its
faces, and any source acting inside it. A \emph{flux} is the amount transported
across a surface per unit area and unit time. Multiplying by face area gives
a transport rate. Here a face length times velocity gives volume flow rate
per unit depth.

For the incompressible cavity, the mass inside a fixed cell is constant.
Its continuity balance is therefore
\[
\boxed{\text{total mass flow rate out of a cell}
       -\text{total mass flow rate into it}=0.}
\]
Momentum has a related balance, with acceleration, pressure forces and viscous
transport as well as momentum carried by fluid crossing the faces.

\textbf{How the equations become numerical.}
First integrate a conservation equation over the chosen region. Next express
transport and forces through its faces. Finally approximate the required
face values using a finite set of stored values. The resulting equation is
an algebraic relation between neighboring unknowns. This process is called
\emph{discretization}.

\textbf{Why sharing a face matters.} Fluid leaving one cell through an internal
face enters the neighboring cell through that same face. Using one numerical
face flow rate in both balances makes these internal contributions cancel
when the balances are added. Each individual cell must still balance;
cancellation in the whole domain alone is insufficient.

\newpage
\topic{2 continued}{Finite volume and finite differences}
The finite-difference approach used in Week 1 begins with a differential
equation evaluated at grid points. It replaces derivatives by combinations
of nearby samples. For a smooth function $f(x)$ sampled at equally spaced
positions $x_i$, define $f_i=f(x_i)$ and $h=x_{i+1}-x_i$. For example,
\[
\left.\frac{\mathrm df}{\mathrm dx}\right|_{x_i}
\simeq\frac{f_{i+1}-f_{i-1}}{2h}.
\]
This is a centered approximation to a derivative \emph{at a point}.

The finite-volume approach begins with a conservation balance \emph{over a
region}. Its approximations concern the values and gradients needed at the
region's faces. Those gradients can also be evaluated with differences.
Using a difference formula inside a finite-volume method does not change
the fact that the starting equation is a cell conservation balance.

\begin{center}\begin{tabular}{>{\raggedright\arraybackslash}p{.22\linewidth}>{\raggedright\arraybackslash}p{.33\linewidth}>{\raggedright\arraybackslash}p{.33\linewidth}}\toprule
Question & Finite differences & Finite volume\\\midrule
Starting form & Differential equation at a grid location & Integrated conservation law for a cell\\
Numerical construction & Approximate derivatives from nearby samples & Approximate face transport and cell terms\\
Conservation check & Depends on the chosen discrete construction & Check the balance of every cell using shared face fluxes\\
This course & Week 1 streamfunction--vorticity solver & Week 1.2 pressure--velocity solver\\\bottomrule
\end{tabular}\end{center}
On a uniform Cartesian mesh, the two approaches can produce identical
algebraic formulas for some terms, such as constant-coefficient diffusion.
Finite-difference schemes can also be constructed in conservative form.
The distinction is the construction of the discrete equations, not a claim
that one approach always conserves and the other never can.

For the cavity, the next pages choose where to store each unknown. Then the
continuity integral is evaluated over one cell and the momentum integral over
a region associated with one velocity unknown. Their locations and boundaries
are specified before performing those integrations.
